% SPDX-License-Identifier: BSD-3-Clause
% Diretrizes do projeto VoronoiMeshMaker — versão vigente (P05, 2026-09-28).
% Substitui VoronoiGridMaker/docs/architecture/project_guidelines.tex, mantida como histórico (DEC-026).
% Cada regra cita o requisito (planning/P05_BASELINE.md) e a decisão (planning/DECISIONS.md) de origem.

\section{Diretrizes do projeto VoronoiMeshMaker}
\label{sec:vmm_project_guidelines}

O \texttt{VoronoiMeshMaker} (VMM) é uma biblioteca em \textbf{C++23} que recebe um domínio, as suas regiões e um
conjunto de sítios geradores e entrega um \textbf{modelo de malha de Voronoi pronto para volumes finitos}
(R14, DEC-001). A primeira versão pública é bidimensional e multirregião; a extensão tridimensional usa o mesmo
modelo de dados (DEC-025).

Cada regra abaixo é \textbf{verificável}: a verificação indicada entre colchetes é feita pela integração
contínua (CI) ou pela revisão do João em cada iteração.

\subsection{Princípios de arquitetura}

\begin{enumerate}
  \item \textbf{Sem funções virtuais, sem herança e sem enum de despacho.} A extensão é feita por templates,
        concepts, traits, policies, composição e registros abertos. (R3, DEC-024)
        [CI: busca de \texttt{virtual} e de listas de classes base nos fontes da biblioteca; revisão.]
  \item \textbf{SOLID por polimorfismo estático.} Responsabilidade única por módulo; extensão sem modificar o
        núcleo (novas policies e novos registros); contratos expressos como concepts; dependência de concepts, não
        de tipos concretos. (R4)
        [Revisão; concepts documentados no P06.]
  \item \textbf{Orientação a dados (DOD).} Dados em estruturas de arrays contíguos, com identificadores fortes e
        conectividade compacta (tipo CSR); dados base separados de dados derivados. (R4)
        [Revisão do modelo de dados no P06.]
  \item \textbf{Separação de responsabilidades (SoC).} O grafo de dependências entre módulos não tem ciclos; o
        núcleo não depende de IO nem do backend geométrico. (R4)
        [CI: script que extrai os includes entre módulos e falha em ciclo ou em dependência proibida.]
  \item \textbf{Modelo de domínio (DDD).} Os conceitos do domínio (meio, região, interface, patch, sítio,
        célula, face) têm tipos próprios, com os nomes usados na documentação. (R5, R6)
        [Revisão.]
  \item \textbf{API pública pequena, estável e só com tipos do VMM.} Nenhum tipo do CGAL, do PETSc ou de outra
        biblioteca externa aparece na API. (R8, R19, R20)
        [CI: headers públicos compilam sem o CGAL no caminho de includes; busca de \texttt{petsc}.]
\end{enumerate}

\subsection{Enums}

\begin{enumerate}
  \item \textbf{Enum de dado é permitido}: valor que descreve um estado ou uma categoria e não escolhe
        implementação (por exemplo, severidade de um erro, idioma, lado de uma face). (DEC-024)
  \item \textbf{Enum de despacho é proibido}: enum cujo valor, num \texttt{switch} ou cadeia de \texttt{if},
        seleciona entre implementações alternativas. Essa escolha é feita por policy (tempo de compilação) ou por
        registro aberto de callables (tempo de execução). (R3, DEC-024)
        [Revisão; o prompt REV verifica cada enum novo.]
\end{enumerate}

\subsection{Domínio, regiões e sítios}

\begin{enumerate}
  \item O domínio tem $N \geq 1$ regiões com fronteira fixa. Cada região pertence a um \textbf{meio} de um
        registro aberto; um meio tem de zero a várias regiões. (R5)
  \item As regiões são declaradas por \textbf{precedência}: cada região declarada depois cobre as anteriores. A
        declaração é convertida numa \textbf{partição explícita}, verificada por um validador que rejeita vazios
        (salvo região de fundo), informa regiões anuladas pela ordem e informa lascas abaixo de uma fração da
        escala local. (R15, DEC-018)
        [GTest do validador.]
  \item Interfaces entre regiões são \textbf{conformes} e rotuladas pelo par de regiões, inclusive entre regiões
        do mesmo meio. (R6)
  \item Os sítios fornecidos pelo usuário ficam estritamente no interior das regiões. Sítios auxiliares (por
        exemplo, espelhados) são internos à biblioteca e não aparecem na malha entregue.
  \item Os sítios são os pontos de referência dos volumes finitos. A entidade computacional é a célula de
        Voronoi \textbf{recortada} pelo domínio e pelas interfaces.
  \item A malha é \textbf{estática}: depois de construída, só expõe acesso \texttt{const}. (R7)
\end{enumerate}

\subsection{Invariantes, tolerâncias e determinismo}

\begin{enumerate}
  \item Toda malha entregue satisfaz os \textbf{invariantes}: a soma das medidas das células é igual à medida
        do domínio e à de cada região; os vetores de área de cada célula somam zero; owner e neighbour são
        consistentes; toda face interna liga exatamente duas células; toda face de contorno liga exatamente uma;
        toda face de interface tem uma célula de cada região. (R6, R16, DEC-011)
        [Verificador de invariantes em todos os testes de integração e no benchmark.]
  \item Tolerâncias são \textbf{relativas} à diagonal $L$ da caixa envolvente do domínio; o padrão é
        $10^{-12}L$. Constantes absolutas de tolerância são proibidas no código geométrico. (R17, DEC-020)
        [CI: busca de constantes absolutas; testes com o domínio escalado por $10^{-3}$ e por $10^{6}$.]
  \item Predicados geométricos são exatos (backend CGAL). A orientação positiva é a do CGAL; entidades
        orientadas ao contrário são reorientadas de forma determinística.
  \item \textbf{Determinismo}: mesma entrada, configuração, semente e build produzem malha idêntica bit a bit,
        incluindo a numeração. Outra plataforma ou outra ordem de inserção dos sítios produzem a mesma topologia
        e a mesma numeração canônica, com geometria dentro da tolerância. Geradores aleatórios usam distribuições
        próprias e portáveis, não as de \texttt{<random>}. (R18, DEC-020)
        [Teste de repetição com soma de verificação; teste com permutações dos sítios.]
\end{enumerate}

\subsection{Backend geométrico}

\begin{enumerate}
  \item O CGAL é o único backend geométrico e fica atrás de uma policy interna; ele responde a perguntas
        geométricas, e o VMM constrói e possui a malha. (DEC-002, DEC-005, DEC-006, DEC-007)
  \item Os headers do CGAL só são incluídos nos fontes do alvo \texttt{vmm\_backend\_cgal}; o alvo
        \texttt{vmm\_core} não depende do CGAL. (R19, DEC-007)
        [CI: teste ``nenhum header público inclui o CGAL''.]
  \item O tipo escalar público é \texttt{vmm::Real = double}. (R20)
\end{enumerate}

\subsection{Erros e mensagens}

\begin{enumerate}
  \item O VMM tem um subsistema próprio de erros, sem herança nem funções virtuais, independente do CGAL, do
        PETSc e das aplicações. (R10, DEC-016)
  \item Falhas previsíveis e recuperáveis são devolvidas por \texttt{std::expected}. Condições excepcionais usam
        uma exceção própria do VMM, que \textbf{não deriva} de \texttt{std::exception}; a documentação avisa que
        \texttt{catch (const std::exception\&)} não a captura. (R11, DEC-016)
  \item Cada erro tem código estável, categoria, mensagem, origem, contexto, \texttt{std::source\_location} e,
        quando existir, a entidade relacionada (\texttt{SiteId}, \texttt{CellId}, \texttt{FaceId},
        \texttt{RegionId}, \texttt{PatchId}). (R10)
  \item A lógica do erro é separada do texto; as mensagens vêm de um catálogo em português e em inglês, com o
        idioma escolhido por configuração explícita. (R10)
        [Teste de catálogo: todo código tem texto nos dois idiomas.]
  \item O logging passa por uma fachada própria; nenhuma biblioteca de logging externa entra no núcleo. (DEC-009)
\end{enumerate}

\subsection{Linguagem, compiladores e nomes}

\begin{enumerate}
  \item Padrão \textbf{C++23}, com \texttt{std::expected}. Compiladores mínimos: \textbf{GCC 14} e
        \textbf{Clang 18}. Recursos com suporte incompleto nesses compiladores (por exemplo, \texttt{std::mdspan})
        não são usados. (R12, DEC-014, DEC-022)
        [CI com as versões mínimas e as mais recentes.]
  \item Nome do projeto e da biblioteca: \texttt{VoronoiMeshMaker}. Namespace público: \texttt{vmm}.
        (R13, DEC-014, DEC-022)
  \item Includes agrupados na ordem: biblioteca padrão, bibliotecas externas, VMM; cada grupo em ordem
        alfabética, com os comentários separadores do projeto. (AGENTS.md)
  \item Identificadores, comentários técnicos e documentação Doxygen em inglês; documentação de usuário e
        entregáveis de planejamento em português, com tradução para o inglês na documentação pública. (R29)
\end{enumerate}

\subsection{Dependências, versões e licença}

\begin{enumerate}
  \item Obrigatórias: CGAL e Boost (headers). Multiprecisão padrão: Boost.Multiprecision. Opcionais: GMP, MPFR,
        TBB. GoogleTest só nos testes. Qualquer outra dependência exige nova DEC. (R21, DEC-004, DEC-009)
        [CI compara os \texttt{find\_package} com a lista permitida.]
  \item Acompanham-se as versões \textbf{estáveis} mais recentes. O CMake declara a versão mínima; a CI mantém as
        versões testadas; cada build, teste e relatório registra as versões usadas. Betas, release candidates e
        builds internos não são usados. (R22, DEC-012)
  \item \texttt{vmm\_core} sob BSD-3-Clause; \texttt{vmm\_backend\_cgal} sujeito às obrigações da GPL dos pacotes
        do CGAL que usa; todo arquivo-fonte com cabeçalho SPDX. Isso não é aconselhamento jurídico. (R23, DEC-008)
        [CI verifica o SPDX.]
  \item Build sem \texttt{-ffast-math}; \texttt{VMM\_ENABLE\_NATIVE\_ARCH} e \texttt{VMM\_ENABLE\_LTO} desligados
        por padrão; saída de exemplos fora da árvore de fontes; nenhum binário versionado; fins de linha LF.
        (R24, DEC-012)
        [CI com \texttt{-Werror} nos alvos do VMM; \texttt{git status} limpo depois do build.]
\end{enumerate}

\subsection{Testes}

\begin{enumerate}
  \item Toda classe \textbf{não trivial} tem o seu \texttt{ut\_<Classe>.cpp}; classe modificada tem o teste
        atualizado. Classe trivial é um agregado sem invariantes nem comportamento (tag, opção sem validação,
        registro de dados); a lista de classes triviais é declarada e revisada. (R1, DEC-023)
  \item \texttt{tests/} espelha a árvore da biblioteca, com descoberta automática. (R2)
  \item Cobertura do código novo: pelo menos 90\% das linhas e 80\% dos ramos por arquivo; exceções por arquivo,
        com justificativa. (R25, DEC-023)
        [Relatório do gcovr na CI.]
  \item Além dos testes por classe: testes de propriedade (invariantes), golden files com tolerância relativa e
        casos patológicos (sítios quase cocirculares, colados à interface, regiões finas, escalas extremas)
        desde a primeira entrega. (R26, DEC-011)
  \item Primeiro os testes das funções públicas de cada classe, depois os testes de integração entre classes;
        só então um exemplo novo entra na galeria. (R30, AGENTS.md)
\end{enumerate}

\subsection{Qualidade e desempenho}

\begin{enumerate}
  \item A biblioteca calcula e relata métricas de qualidade para volumes finitos (não ortogonalidade, distorção,
        faces e células mínimas, razão de aspecto); os limites de aceitação são configuráveis. (R14)
  \item Benchmark leve (casos B1, B2 e B3) a cada entrega, com tempo por fase e pico de memória; as metas de
        tempo e memória podem ser recalibradas uma única vez com a linha de base do P07. (R27, DEC-020, DEC-021)
\end{enumerate}

\subsection{Conectividade, reordenação e saída}

\begin{enumerate}
  \item A biblioteca fornece views para células e faces internas, de contorno (por patch) e de interface (por par
        de regiões), e a adjacência em forma compacta, suficiente para montar matrizes esparsas sem reconstruir a
        geometria. (R9, DEC-015)
  \item A reordenação dos volumes guarda explicitamente a permutação e a permutação inversa.
  \item A visualização (VTK) é separada da persistência (formato nativo). Saída da 0.1: formato nativo
        (prioritário), VTK XML e OpenFOAM polyMesh; os escritores ficam fora do núcleo. (R28, DEC-019)
\end{enumerate}

\subsection{Documentação}

\begin{enumerate}
  \item Referência da API no modelo do PETSc (Synopsis, Parameters, Notes, Level, See Also, Location, Examples)
        e galeria sphinx-gallery com a paleta ``Estuário''; tema PyData. (R29)
  \item Um exemplo que falha quebra o build da documentação. (R29)
\end{enumerate}

\subsection{Ambiente local e integração contínua}

\begin{enumerate}
  \item O \texttt{AGENTS.md} rege o trabalho dos agentes na árvore local do João: ambiente único (WSL
        \texttt{Ubuntu-26.04-Test}), sem commit nem push, e resultados de outro ambiente não valem como evidência
        para a árvore local.
  \item A CI pública valida os commits publicados. As duas evidências são complementares e não se substituem.
\end{enumerate}